% Compact two-page academic CV. Publications are generated from the site data.
\documentclass[letterpaper,10pt]{article}
\usepackage[margin=0.6in,footskip=0.25in]{geometry}
\usepackage[T1]{fontenc}
\usepackage[scaled=0.95]{helvet}
\renewcommand{\familydefault}{\sfdefault}
\usepackage{microtype,xcolor,hyperref,tabularx,array,enumitem,fancyhdr}
\definecolor{Graphite}{HTML}{1D1D1F}
\definecolor{Muted}{HTML}{626269}
\definecolor{Accent}{HTML}{5964A0}
\definecolor{Rule}{HTML}{DDDEE3}
\color{Graphite}
\hypersetup{colorlinks=true,urlcolor=Graphite,linkcolor=Graphite,pdftitle={Chunghyun Park - Curriculum Vitae},pdfauthor={Chunghyun Park}}
\setlength{\parindent}{0pt}
\setlength{\parskip}{0pt}
\setlength{\tabcolsep}{0pt}
\setlist[itemize]{leftmargin=12pt,itemsep=1pt,topsep=3pt,parsep=0pt}
\pagestyle{fancy}\fancyhf{}
\renewcommand{\headrulewidth}{0pt}
\fancyfoot[L]{\fontsize{8}{10}\selectfont\color{Muted}Chunghyun Park \quad | \quad Curriculum vitae}
\fancyfoot[R]{\fontsize{8}{10}\selectfont\color{Muted}Last updated: \today\quad | \quad\thepage}
\newcommand{\cvsection}[1]{\par\addvspace{12pt}{\fontsize{11}{13}\selectfont\bfseries #1}\par\vspace{4pt}{\color{Rule}\hrule}\vspace{6pt}}
\newcommand{\entry}[3]{\textbf{#1}\hfill{\small\color{Muted}#2}\par #3\par\vspace{5pt}}
\newcommand{\dated}[2]{\par\hangindent=38pt\hangafter=1\noindent\makebox[38pt][l]{\color{Muted}#1}#2\par\vspace{3pt}}
\newcommand{\contribmark}[1]{\raisebox{0.45ex}{\color{Muted}\fontsize{8.5}{8.5}\selectfont\bfseries #1}}
\newcommand{\publication}[4]{\begin{minipage}{\linewidth}\fontsize{9.2}{11.2}\selectfont
\makebox[27pt][l]{\color{Muted}\textbf{#1}}\parbox[t]{\dimexpr\linewidth-27pt\relax}{\textbf{#2}\par {\color{Muted}#3}\par #4}
\end{minipage}\par\vspace{6pt}}
\begin{document}
{\fontsize{29}{33}\selectfont\bfseries Chunghyun Park}\par\vspace{5pt}
{\color{Muted}Postdoctoral researcher, POSTECH Institute of AI}\par\vspace{5pt}
{\small\href{mailto:p0125ch@postech.ac.kr}{p0125ch@postech.ac.kr}\quad\textperiodcentered\quad
\href{https://chrockey.github.io}{chrockey.github.io}\quad\textperiodcentered\quad
\href{https://github.com/chrockey}{GitHub}\quad\textperiodcentered\quad
\href{https://scholar.google.com/citations?user=5ABvjQcAAAAJ}{Google Scholar}}

\cvsection{Research}
3D perception and dexterous manipulation: understanding 3D environments and anticipating the effects of robot actions, with work spanning geometric learning, open-vocabulary and active perception, and world action modeling

\cvsection{Experience}
\entry{POSTECH Institute of AI \textnormal{\,|\, Postdoctoral researcher}}{Mar. 2026 -- Present}{
3D perception and dexterous manipulation\par
{\small\color{Muted}Advisor: Prof.\ Minsu Cho \textperiodcentered\ Pohang, Republic of Korea \textperiodcentered\ Mandatory military service through Aug. 2027}}
\entry{NVIDIA \textnormal{\,|\, Research collaborator}}{Mar. 2025 -- Feb. 2026}{
Open-vocabulary 3D instance segmentation: SpaCeFormer (ICML 2026)\par
{\small\color{Muted}Collaborators: Chris Choy and Jan Kautz \textperiodcentered\ Santa Clara, USA (remote)}}
\entry{NVIDIA \textnormal{\,|\, Research intern}}{Dec. 2023 -- Feb. 2025}{
Foundation data and models for 3D segmentation: Mosaic3D (CVPR 2025)\par
{\small\color{Muted}Mentors: Jaesung Choe and Yu-Chiang Frank Wang \textperiodcentered\ Taipei, Taiwan (remote)}}

\cvsection{Education}
\textbf{Pohang University of Science and Technology (POSTECH)}\par\vspace{4pt}
\entry{Ph.D., Artificial Intelligence \textnormal{\small\color{Muted}\,|\, Advisor: Prof.\ Minsu Cho}}{2022 -- 2026}{
{\small Thesis: Learning 3D Visual Perception for Geometric and Semantic Generalization}}
\entry{M.S., Artificial Intelligence \textnormal{\small\color{Muted}\,|\, Advisor: Prof.\ Jaesik Park}}{2020 -- 2022}{
{\small Thesis: Fast Point Transformer for Large-scale 3D Scene Understanding}}
\textbf{B.S., Mechanical Engineering}\hfill{\small\color{Muted}2014 -- 2019}\par

\cvsection{Honors}
\dated{2026}{CVPR Doctoral Consortium (24 selected graduates worldwide)}
\dated{2026}{Outstanding Presentation Award, KCCV Doctoral Colloquium}
\dated{2026}{KCCV Doctoral Colloquium (9 selected graduates in South Korea)}
\dated{2026}{Bronze Prize, IPIU Best Paper Award}
\dated{2025}{POSTECHIAN Fellowship (Innovation)}
\dated{2024}{CVPR Outstanding Reviewer (top 2\% of 9,872 reviewers)}
\dated{2023}{Grand Prize, BK21 Best Paper Award}
\dated{2022}{Qualcomm Innovation Fellowship}
\dated{2022}{Silver Prize, Samsung HumanTech Paper Award}

\cvsection{Talks}
\dated{2026}{\textit{Embodied Vision for Real-World Perception and Action}, KCCV}
\dated{2026}{\textit{Generalizable 3D Perception for Physical AI}, Chung-Ang University}
\dated{2024}{\textit{Equivariant Networks for 3D Vision}, POSTECH CSED703J}
\dated{2023}{\textit{Stable and Consistent Prediction of 3D Characteristic Orientation}, Korea AI Summit}
\cvsection{Professional service}
\textbf{Reviewer}\quad TPAMI, IJCV; CVPR, ICCV, ECCV, NeurIPS, ICLR

\newpage
\cvsection{Publications}
{\small\color{Muted}\contribmark{\textasteriskcentered}, \contribmark{\textdagger}: Equal contribution within each symbol group}\par\vspace{6pt}
\input{publications}

\cvsection{Patent}
{\small\textbf{Hardware Apparatus and Method for Predicting Structural Characteristic Orientation of Point Cloud of Structural Object}\par
U.S. Patent Application No. US 2025/0191218 A1, 2025}
\end{document}
